% 数学ノート集・共通テンプレート v2（積分論2026年度第1部の版面を参考）
% A5判。章内で定義・定理・例・注意等の番号を共通化。
% コンパイル： platex -interaction=nonstopmode math_note_template_v2.tex (2回)
%              dvipdfmx math_note_template_v2.dvi
\documentclass[a5paper,10pt,onecolumn]{jsarticle}
\usepackage[dvipdfmx,top=22mm,bottom=20mm,inner=19mm,outer=18mm,headheight=13pt,headsep=6mm]{geometry}
\usepackage{amsmath,amssymb,amsthm,mathtools}
\usepackage{enumitem,fancyhdr,titlesec}
\usepackage[dvipdfmx,bookmarks=true,colorlinks=false]{hyperref}
\usepackage{pxjahyper}

% ----- 差し替える文書情報 -----
\newcommand{\NoteTitle}{積位相・商位相と普遍性}
\newcommand{\NoteSeries}{数学ノート集}

% ----- 文字・組版 -----
% 参考PDFの質素な明朝＋Computer Modernの印象を引き継ぎながら，
% A5判で窮屈にならないよう本文の行間を少し確保する．
\linespread{1.13}
\setlength{\parindent}{1zw}
\setlength{\parskip}{0pt}
\setlength{\abovedisplayskip}{6pt plus 2pt minus 2pt}
\setlength{\belowdisplayskip}{6pt plus 2pt minus 2pt}
\setlength{\abovedisplayshortskip}{3pt plus 2pt minus 1pt}
\setlength{\belowdisplayshortskip}{4pt plus 2pt minus 1pt}
\setlist[enumerate]{topsep=3pt,itemsep=1pt,parsep=0pt,leftmargin=2.4zw}
\titleformat{\section}{\normalfont\normalsize\bfseries}{\thesection}{0.8zw}{}
\titlespacing*{\section}{0pt}{2.0ex plus .4ex}{.9ex}
\titleformat{\subsection}{\normalfont\normalsize\bfseries}{\thesubsection}{0.7zw}{}
\titlespacing*{\subsection}{0pt}{1.5ex plus .3ex}{.6ex}
\raggedbottom

% ----- 柱とページ番号：ページ上部、罫線なし -----
\fancypagestyle{mathnote}{%
  \fancyhf{}%
  \fancyhead[L]{\footnotesize \NoteSeries\ （\NoteTitle）}%
  \fancyhead[R]{\footnotesize \thepage}%
  \renewcommand{\headrulewidth}{0pt}%
  \renewcommand{\footrulewidth}{0pt}%
}
\pagestyle{mathnote}

% ----- 節内通し番号の定義・定理群 -----
\newtheoremstyle{mathnotetheorem}
 {4.5pt}{4.5pt}{\normalfont}{}%
 {\bfseries}{}{0.65zw}%
 {\thmname{#1}\thmnumber{ #2}\thmnote{（#3）}}
\theoremstyle{mathnotetheorem}
\newtheorem{theorem}{定理}[section]
\newtheorem{definition}[theorem]{定義}
\newtheorem{lemma}[theorem]{補題}
\newtheorem{proposition}[theorem]{命題}
\newtheorem{corollary}[theorem]{系}
\newtheorem{example}[theorem]{例}
\newtheorem{remark}[theorem]{注意}
\renewcommand{\proofname}{証明}
\renewcommand{\qedsymbol}{$\square$}
\makeatletter
\renewenvironment{proof}[1][\proofname]{%
  \par\pushQED{\qed}\normalfont\topsep4\p@\@plus2\p@\relax
  \trivlist\item[\hskip\labelsep\bfseries #1]\ignorespaces
}{\popQED\endtrivlist\@endpefalse}
\makeatother

% ----- 使うノートで適宜拡張する記号 -----
\newcommand{\R}{\mathbb{R}}
\newcommand{\norm}[1]{\lVert #1\rVert}
\newcommand{\abs}[1]{\lvert #1\rvert}

\begin{document}
\section*{読了後にできること}
本稿を読むための前提は，積位相・商位相の定義と普遍性．
本稿の具体的な到達課題は，商空間の位相を具体的に記述し，連続性を商写像で検証する．


\section{始位相と終位相}
本稿では $A\subset B$ は $A=B$ を許す包含を表し，
$f^{-1}(U)$ は写像 $f$ による集合 $U$ の逆像を表す．
添字集合は特に断らない限り任意の集合とし，有限性を要求しない．
\begin{definition}
集合 $X$ と位相空間族 $(Y_i)_{i\in I}$ および写像
$f_i:X\to Y_i$ に対し，
全ての $f_i^{-1}(U)$（$U$ は $Y_i$ の開集合）
で生成される $X$ の位相を始位相という．
逆に位相空間族 $(X_i)$ と集合 $Y$ および
$g_i:X_i\to Y$ に対し，
\[
 \mathcal T_Y=\{V\subset Y:g_i^{-1}(V)\text{ は全ての }i
 \text{ で開}\}
\]
を終位相という．
\end{definition}
\begin{proposition}[始位相の普遍性]\label{prop:initial}
$X$ に始位相を入れると，
任意の位相空間 $Z$ と写像 $h:Z\to X$ に対して
\[
 h\text{ は連続}\quad\Longleftrightarrow\quad
 f_i\circ h\text{ は全て連続}
\]
が成り立つ．
\end{proposition}
\begin{proof}
必要性は $f_i$ が始位相の定義で連続となることと合成の連続性による．
逆に，各 $f_i\circ h$ が連続なら
$h^{-1}(f_i^{-1}(U))=(f_i\circ h)^{-1}(U)$ は開である．
始位相の準基底の逆像が開なので $h$ は連続である．
\end{proof}
\begin{proposition}[終位相の普遍性]\label{prop:final}
$Y$ に終位相を入れると，
任意の位相空間 $Z$ と写像 $h:Y\to Z$ に対して
\[
 h\text{ は連続}\quad\Longleftrightarrow\quad
 h\circ g_i\text{ は全て連続}
\]
となる．
\end{proposition}
\begin{proof}
必要性は合成の連続性による．
逆に開集合 $W\subset Z$ について
$g_i^{-1}(h^{-1}(W))=(h\circ g_i)^{-1}(W)$ は全て開だから，
終位相の定義から $h^{-1}(W)$ は開である．
\end{proof}
\section{部分空間・直積・直和}
\begin{definition}
$A\subset X$ に対し
$\{A\cap U:U\subset X\text{ 開}\}$ を相対位相という．
位相空間族 $(X_i)$ の集合としての積
$\prod_iX_i$ に，座標射影 $\pi_i$ に関する始位相を入れたものを
積空間という．
\end{definition}
\begin{lemma}[積位相の基底]\label{lem:productbase}
$\prod_iX_i$ の積位相は
\[
 \prod_iU_i,\qquad
 U_i\subset X_i\text{ 開},\quad U_i=X_i\text{ を除く添字は有限}
\]
を基底とする．
\end{lemma}
\begin{proof}
始位相の準基底は $\pi_i^{-1}(U_i)$ である．
有限個の準基底元を交叉すると表示した形となり，
逆に表示した形は有限個の準基底元の交叉である．
\end{proof}
\begin{proposition}[直積の普遍性]
$f:Z\to\prod_iX_i$ は，
全ての座標 $\pi_i\circ f$ が連続であるとき，かつそのときに限り連続である．
\end{proposition}
\begin{proof}
積位相は射影に関する始位相なので
命題\ref{prop:initial}を適用する．
\end{proof}
\begin{definition}
集合としての互いに素な和 $\coprod_iX_i$ に，
$U$ が開であることと各 $X_i$ への制限 $U\cap X_i$ が開であることが
同値となる位相を入れたものを位相的直和という．
\end{definition}
\begin{proposition}[直和の普遍性]
$f:\coprod_iX_i\to Z$ が連続であることと，
各包含写像 $\iota_i:X_i\to\coprod_iX_i$ との合成
$f\circ\iota_i$ が連続であることは同値である．
\end{proposition}
\begin{proof}
直和の位相は各 $\iota_i$ に関する終位相であるから
命題\ref{prop:final}より従う．
\end{proof}
\section{商位相と貼り合わせ}
\begin{definition}
全射 $q:X\to Y$ に対し
$V\subset Y$ が開であることと $q^{-1}(V)$ が開であることが
同値となる位相を $Y$ に入れる．これを商位相といい，
$q$ を商写像という．
\end{definition}
\begin{theorem}[商写像による降下]\label{thm:quotient}
$q:X\to Y$ を商写像とし，$f:X\to Z$ を連続写像とする．
$f$ が $q$ の各ファイバー上定数であることと，
連続写像 $\bar f:Y\to Z$ が存在して
$f=\bar f\circ q$ となることは同値である．
その場合 $\bar f$ は一意である．
\end{theorem}
\begin{proof}
$f$ が各ファイバー上定数なら
$\bar f(q(x))=f(x)$ と置いてwell-definedな写像を得る．
$Z$ の開集合 $W$ に対して
$q^{-1}(\bar f^{-1}(W))=f^{-1}(W)$ は開なので，
商位相の定義から $\bar f^{-1}(W)$ は開である．
全射性よりこの写像は一意である．
逆に因子分解があれば同じファイバー上の二点は同じ像を持つ．
\end{proof}
\begin{lemma}[開全射・閉全射は商写像]
連続全射 $q:X\to Y$ が開写像または閉写像なら商写像である．
\end{lemma}
\begin{proof}
$V\subset Y$ が $q^{-1}(V)$ を開とすると，
$q$ が開写像なら $V=q(q^{-1}(V))$ は開である．
閉写像の場合は
$Y\setminus V=q(X\setminus q^{-1}(V))$ が閉なので $V$ は開である．
\end{proof}
\begin{theorem}[貼り合わせ補題]\label{thm:pasting}
$X=\bigcup_{i\in I}U_i$ を開被覆とし，
連続写像 $f_i:U_i\to Y$ が重なりで一致するとする．
このとき一意の連続写像 $f:X\to Y$ が存在して
$f|_{U_i}=f_i$ を満たす．
有限個の閉集合 $F_1,\dots,F_m$ による被覆でも同じ結論が成立する．
\end{theorem}
\begin{proof}
写像 $f$ は点ごとに $f(x)=f_i(x)$ と置けば一致条件から定まる．
開被覆なら $f^{-1}(V)\cap U_i=f_i^{-1}(V)$ が
$U_i$ の開集合なので $X$ でも開であり，
$f^{-1}(V)=\bigcup_i(f^{-1}(V)\cap U_i)$ は開となる．
有限閉被覆なら閉集合 $C\subset Y$ に対して
$f^{-1}(C)\cap F_i=f_i^{-1}(C)$ は $F_i$ で閉，
従って $X$ で閉である．
有限和は閉であるから $f^{-1}(C)$ は閉となる．
\end{proof}
\section{例：円周の商表示}
\begin{proposition}
$[0,1]$ の両端を同一視した商空間は単位円周
$S^1=\{z\in\mathbb C:|z|=1\}$ と同相である．
\end{proposition}
\begin{proof}
$f:[0,1]\to S^1$，$f(t)=e^{2\pi it}$ は連続全射であり，
$f(s)=f(t)$ は $s=t$ または $\{s,t\}=\{0,1\}$
と同値である．よって商集合への全単射を誘導する．
区間はコンパクトで円周はHausdorffなので，
連続全単射は同相である：閉集合の像はコンパクトゆえ閉となり，
逆写像が連続となるためである．
\end{proof}

\section{到達課題の完全解答}
\noindent\textbf{区間の端点の同一視から円周を得る．}
\begin{proposition}[端点の貼り合わせ]
$[0,1]$ の $0\sim1$ を同一視して得る商空間
$Q=[0,1]/\sim$ は円周 $S^1$ と同相である．
\end{proposition}
\begin{proof}
$f:[0,1]\to S^1$ を $f(t)=e^{2\pi it}$ と置く．
$f(0)=f(1)$ だから商写像
$q:[0,1]\to Q$ の普遍性により
$\bar f:Q\to S^1$ が存在して $f=\bar f\circ q$ となり，
$\bar f$ は連続である．
指数関数の周期性から $\bar f$ は全単射である．
$Q$ はコンパクト，$S^1$ はHausdorffなので
連続全単射 $\bar f$ は同相写像である．
\end{proof}
\begin{remark}
商位相を定めるだけでは，商空間が図形として
何に対応するかはまだ判明しない．
具体的な連続全単射とコンパクト性・Hausdorff性を
組み合わせることが重要である．
\end{remark}
\noindent\textbf{到達確認．}
$[0,2\pi]$ の両端を同一視した商空間が
$S^1$ と同相であることを証明せよ．
\par\noindent\textbf{解答．}
$t\mapsto e^{it}$ は連続で両端で一致するため
商写像を通じて連続全単射を誘導する．
定義域の商空間はコンパクトで値域はHausdorffなので同相である．
\end{document}
